%ÚVODNÍ BALÍKY, VZHLED STRÁNKY
\PassOptionsToPackage{dvipsnames}{xcolor} % Barvičky musí být předané před startem beamer-u
\documentclass[aspectratio=169]{beamer} %Poměr stran 16:9,
\usetheme{Szeged}
% Možné argumenty (doporučuji): Metropolis
% Možné argumenty (další): AnnArbor, Antibes, Bergen, Berkeley, Berlin, Boadilla, boxes, CambridgeUS, Copenhagen, Darmstadt, default, Dresden, Frankfurt, Goettingen, Hannover, Ilmenau, JuanLesPins, Luebeck, Madrid, Malmoe, Marburg, Montpellier, PaloAlto, Pittsburgh, Rochester, Singapore, Szeged, Warsaw
\usecolortheme{beaver}
% Možné argumenty: albatross, beaver, beetle, crane, dolphin, dove, fly, lily, orchid, rose, seagull, seahorse, whale, wolverine
% https://deic.uab.es/~iblanes/beamer_gallery/individual/CambridgeUS-wolverine-structuresmallcapsserif.html

\setbeamertemplate{navigation symbols}{} %Vymaže šipky na slidech
%\setbeamercovered{transparent} %Ukazuje body, které teprve přijdou zašedlé

%Zde jde odkomentovat, když chci prezentaci vytisknout na papír, např. abych si udělal poznámky tužkou. Taky jde "2 on 1" varianta.
%\usepackage{pgfpages}
%\pgfpagesuselayout{4 on 1}[a4paper, border shrink=5mm]

% Různé vychytávky
\usepackage{multicol} % Sloupce
\usepackage[shortlabels]{enumitem} % Zkrácené číslování
% Grafiky - vkládání obrázků
\usepackage{graphicx} % Základní
\usepackage{graphbox} % Obalování boxem zvláštních vlastností 1
\usepackage{adjustbox} % Obalování boxem zvláštních vlastností 2
% Matematické zápisy
\usepackage{mathtools} % Rozšíření asmmath - všechna matika.
\usepackage{array}
%\usepackage[shortlabels]{enumitem}

% Čeština a znakové sady
\usepackage[czech, shorthands=off]{babel}
\usepackage[utf8]{inputenc}
\usepackage[T1]{fontenc}
% Moderní a matematické fonty
\usepackage{lmodern}
\usepackage{amsfonts} % Matematické fonty
\usepackage{amssymb} % Matematické symboly
\newcommand{\C}{\ensuremath{\mathbb{C}}}

%Vlastní příkazy
\NewDocumentCommand{\highlight}{ +m }{\textbf{\color{ForestGreen}
#1}}
\NewDocumentCommand{\Highlight}{ +m }{\noindent
\textbf{\color{ForestGreen}
#1}}
\newcommand{\predepsaneKody}[1]{}

%Vlastní barvy
\definecolor{Gray}{gray}{0.5}

%Definice sloupců tabulek typu Centerer, Left a Right
\newcolumntype{C}[1]{>{\centering\let\newline\\\arraybackslash\hspace{0pt}}m{#1}}
\newcolumntype{L}[1]{>{\raggedright\let\newline\\\arraybackslash\hspace{0pt}}m{#1}}
\newcolumntype{R}[1]{>{\raggedleft\let\newline\\\arraybackslash\hspace{0pt}}m{#1}}

\DeclareMathOperator{\tg}{tg}
\DeclareMathOperator{\cotg}{cotg}

% ===  Krajíčkovské úpravy  ===
\usepackage{xparse}
\usepackage{parskip}

% Matematický font
\usefonttheme[onlymath]{serif}

% Odkazy
\usepackage{hyperref}
\newcommand{\link}[2]{\href{#1}{\color{Cyan}
\underline{\smash{#2}}}}

% PGF Grafy
\usepackage{pgfplots}
\pgfplotsset{compat=1.18}
\pgfplotsset{height=200pt} %velikost
\usetikzlibrary{calc}

% Lepší tabulky.
\usepackage{multirow}
\renewcommand{\arraystretch}{1.5}

% Škrtání v text
\usepackage{cancel}
\renewcommand{\CancelColor}{\color{Red}}

% Animace
\usepackage{animate}

% === Úpravy pro TUTO prezentaci ===
\newcommand{\reseni}[1]{\hfill \hyperlink{#1}{\beamergotobutton{\strut Řešení}}}
\newcommand{\zpet}[1]{\hyperlink{#1}{\beamerreturnbutton{\strut Zpět}} \hfill }
\setbeamercolor{button}{bg=BrickRed}
\setbeamercolor{block title alerted}{fg=BrickRed, bg=BrickRed!20}
\setbeamercolor{block body alerted}{fg=black, bg=BrickRed!10}
\setbeamertemplate{blocks}[rounded][shadow=true]

% ===   ÚVODNÍ SLIDE   ===
\title{Analytická geometrie}
\author{II. ročník}
\date{Aktualizace: \today}


\begin{document}
  \maketitle

  % JÁDRO PREZENTACE
  \section{Úvod - objekty}

  \begin{frame}{K čemu to všechno?}
    Analytická geometrie (AnaGeo) je jedním ze způsobů, jak nahlížet na
    geometrické problémy (např. kde se protnou přímky, protnou se vůbec? Co je
    to přímka?)
	
	\medskip
    AnaGeo používá nástroje analýzy (rovnice, vlastnosti rovnic a funkcí) - je třeba
    mít na paměti, co rovnice a čísla reprezentují, protože se na první pohled bude
    zdát, že jsme se vrátili do 1. ročníku 
	
	\medskip
    Nicméně vše zasadíme do \textbf{vektorového prostoru} 
	
	\medskip
    \textbf{Příklad:} Bod ležící na přímce
  \end{frame}

  \begin{frame}{Pojmy}
    K tomu, abychom mohli bezpečně pokračovat, je třeba si udělat společnou slovní
    zásobu (kterou budu aktivně využívat):
	
	\medskip
    \begin{tabular}{C{3cm} C{3cm} C{3cm} C{3cm}}
      Vektor            & Skalár           & Přímka           & Směrnice    \\
      Bod               & Směrový vektor   & Normálový vektor & Obecný tvar \\
      Parametrický tvar & Parametr         & Směrnicový tvar  & Rovina      \\
      Skalární součin   & Vektorový součin & Posun o vektor   & Odchylka    \\
      Velikost vektoru  & Nulový vektor    & Počátek          & ...         \\
    \end{tabular}
  \end{frame}

  \begin{frame}{Bod x Vektor}
    \begin{alertblock}{Bod}
      Bodem rozumíme prvek roviny v kartézské soustavě souřadnic. Značíme velkým
      písmenem a jeho souřadnice ohraničujeme hranatými závorkami:
      \[
        A[2;-3]
      \]
    \end{alertblock}
    Bod pro nás není nic nového - známe jej už do funkcí. Nicméně to co je nové
    je \textbf{vektor}:
    \begin{alertblock}{Vektor}
      Vektorem rozumíme prvek vektorového prostoru. Vektor má \textbf{směr} a \textbf{velikost}.
      Vektor označujeme malým psacím písmenem s šipkou a jeho "směr" značíme do kulatých
      závorek:
      \[
        \vec{v}(2;-3)
      \]
    \end{alertblock}
  \end{frame}

  \begin{frame}{Bod x Vektor}
    Tady obrázky
  \end{frame}

  \begin{frame}{Určování vektorů}
    BUde
  \end{frame}

  \begin{frame}{Posunutí}
    Bude
  \end{frame}

  \begin{frame}[label=C1]{Cvičení - Posun a tvorba vektorů}
    \begin{tabular}{C{3cm} C{3cm} | C{3cm} C{3cm}}
      \multicolumn{2}{c|}{Vytvořte a nakreslete vektory $\vec{v}=B-A$:} & \multicolumn{2}{c}{Posuňte bod $A$ o vektor $\vec{v}$:} \\
      \hline
      $A[2;2]$ $B[-1;2]$                                                & $A[1;1]$ $B[-1;-2]$                                    & $A[1;0]$ $\vec{v}(1,1)$   & $A[1;0]$ $\vec{v}(-1,1)$  \\
      $A[0;2]$ $B[0;-2]$                                                & $A[1;0]$ $B[-1;-2]$                                    & $A[-1;2]$ $\vec{v}(1,-1)$ & $A[2;2]$ $\vec{v}(-1,2)$  \\
      $A[-2;1]$ $B[2;-1]$                                               & $A[1;1]$ $B[2;2]$                                      & $A[-2;2]$ $\vec{v}(2,-1)$ & $A[2;3]$ $\vec{v}(-1,-1)$ \\
      $A[-2;3]$ $B[4;-5]$                                               & $A[0;0]$ $B[3;-2]$                                     & $A[0;2]$ $\vec{v}(0,-1)$  & $A[3;2]$ $\vec{v}(-2,3)$  \\
    \end{tabular}
    \reseni{Ř1}
  \end{frame}

  \begin{frame}{Sčítání vektorů}
    Vektory lze sčítat (jedná se o "sčítání posunů"): \medskip
    Vezmu-li bod $A[2,1]$ a posunu jej nejdříve o vektor $\vec{u}(1,2)$ a poté
    jej posunu o vektor $\vec{v}(2,-1)$, pak je to úplně samé, jako bych jej
    posunul o vektor $\vec{w}(3,1)=\vec{u}+\vec{v}$: \\
    \begin{tikzpicture}[scale=0.7]
      \draw[gray] (-1,-1) grid (6,4);
      \draw[->, thick] (-1.4,0) -- (6.4,0) node[anchor=south east] {$x$};
      \draw[->, thick] (0,-1.4) -- (0,4.4) node[anchor=north west] {$y$};
      \filldraw (2,1) circle (2pt) node[anchor=north east] {$A$};
      \filldraw (5,2) circle (2pt) node[anchor=north west] {$A'$};
      \draw[->, thick, red]
        (2,1) -- (3,3)
        node[pos=0.5, above, sloped] {$\vec{u}$};
      \draw[->, thick, red]
        (3,3) -- (5,2)
        node[pos=0.5, above, sloped] {$\vec{v}$};
    \end{tikzpicture}
    \hspace{1cm}
    \begin{tikzpicture}[scale=0.7]
      \draw[gray] (-1,-1) grid (6,4);
      \draw[->, thick] (-1.4,0) -- (6.4,0) node[anchor=south east] {$x$};
      \draw[->, thick] (0,-1.4) -- (0,4.4) node[anchor=north west] {$y$};
      \filldraw (2,1) circle (2pt) node[anchor=north east] {$A$};
      \filldraw (5,2) circle (2pt) node[anchor=north west] {$A'$};
      \draw[->, thick, red]
        (2,1) -- (3,3)
        node[pos=0.5, above, sloped] {$\vec{u}$};
      \draw[->, thick, red]
        (3,3) -- (5,2)
        node[pos=0.5, above, sloped] {$\vec{v}$};
      \draw[->, thick, red]
        (2,1) -- (5,2)
        node[pos=0.5, below, sloped] {$\vec{u}+\vec{v}$};
    \end{tikzpicture}
  \end{frame}

  \begin{frame}{Střed úsečky}
    \begin{alertblock}{Střed úsečky}
      Střed úsečky $AB$ daná body $A[x_{a},y_{a}]$ a $B[x_{b},y_{b}]$ rozumíme bod
      $S$, pro který platí $|AS|=|SB|$. Souřadnice středu se vypočítají následovně:
      \[
        S[\dfrac{x_a+x_b}{2};\dfrac{y_a+y_b}{2}]
      \]
    \end{alertblock}
    \begin{tikzpicture}[scale=0.7]
      \draw[gray] (-1,0) grid (6,4);
      \draw (2,1) rectangle (5,3);
      \draw[->, thick] (-1.4,0) -- (6.4,0) node[anchor=south east] {$x$};
      \draw[->, thick] (0,-0.4) -- (0,4.4) node[anchor=north west] {$y$};
      \filldraw (2,1) circle (2pt) node[anchor=north east] {$A$};
      \filldraw[red] (3.5,2) circle (2pt) node[above] {$S_{AB}$};
      \filldraw (5,3) circle (2pt) node[anchor=north west] {$B$};
      \draw[->, thick, red]
        (2,1) -- (5,3)
        node[pos=0.7, below, sloped] {$\vec{u}$};
    \end{tikzpicture}
  \end{frame}

  \begin{frame}{Velikost úsečky/vektoru}
    \begin{alertblock}{Střed úsečky}
      Velikost úsečky $AB$ daná body $A[x_{a},y_{a}]$ a $B[x_{b},y_{b}]$ rozumíme
      velikost vektoru $\vec{u}=B-A$. Lze vypočítat následovně:
      \[
        |AB|=\vec{u}=\sqrt{(x_{b}-x_{a})^{2}+(y_{b}+y_{a})^{2}}=\sqrt{u_{x}^{2}+u_{y}^{2}}
      \]
    \end{alertblock}
    \begin{tikzpicture}[scale=0.7]
      \draw[gray] (-1,0) grid (6,4);
      \draw (2,1) -- (5,1) -- (5,3);
      \draw[-latex, thick] (-1.4,0) -- (6.4,0) node[anchor=south east] {$x$};
      \draw[-latex, thick] (0,-0.4) -- (0,4.4) node[anchor=north west] {$y$};
      \filldraw (2,1) circle (2pt) node[anchor=north east] {$A$};
      \filldraw (5,3) circle (2pt) node[anchor=north west] {$B$};
      \draw[->, thick, red]
        (2,1) -- (5,3)
        node[pos=0.7, below, sloped] {$\vec{u}$};
    \end{tikzpicture}
  \end{frame}

  \begin{frame}[label=C2]{Cvičení - Střed úsečky a její velikost}
    \begin{tabular}{C{3cm} C{3cm} | C{3cm} C{3cm}}
      \multicolumn{2}{c|}{Určete středy úseček $AB$ a její velikost} & \multicolumn{2}{c}{Určete velikost vektorů:} \\
      \hline
      $A[2;2]$ $B[-1;2]$                                             & $A[1;1]$ $B[-1;-2]$                         & $\vec{v}(1,1)$  & $\vec{v}(-1,1)$  \\
      $A[0;2]$ $B[0;-2]$                                             & $A[1;0]$ $B[-1;-2]$                         & $\vec{v}(3,-4)$ & $\vec{v}(4,3)$   \\
      $A[-2;1]$ $B[2;-1]$                                            & $A[1;1]$ $B[2;2]$                           & $\vec{v}(2,-1)$ & $\vec{v}(-1,-1)$ \\
      $A[-2;3]$ $B[4;-5]$                                            & $A[0;0]$ $B[3;-2]$                          & $\vec{v}(0,-1)$ & $\vec{v}(-2,3)$  \\
    \end{tabular}
    \reseni{Ř2}
  \end{frame}

  \section{Parametrické vyjádření přímky}

  \begin{frame}{Jak lze vyjádřit přímku?}
    To už jsme měli
  \end{frame}


  \begin{frame}[fragile]{Ukázka}
    \begin{animateinline}
      [autoplay,loop]{10} 
	  	\multiframe{180}{i=-120+1}{ 
			\begin{tikzpicture}[
				scale=0.7,
				line cap=round,
				line join=round,
				grid/.style={gray!35, thin},
				axis/.style={->, thick},
				vector/.style={red, -latex, thick},
				refvector/.style={-latex, thick},
				point/.style={fill=black},
				current point/.style={fill=black},
				vector label/.style={pos=0.5, above, sloped},
				axis label/.style={anchor=north east},
				info label/.style={anchor=west}
			]
				\def\xmin{-5}
				\def\xmax{5}
				\def\ymin{-5}
				\def\ymax{5}
				
				\coordinate (A) at (-2,-2);
				\coordinate (V) at (2,1);
				\pgfmathsetmacro{\t}{-0.025*\i}
				
				\draw[grid] (\xmin,\ymin) grid (\xmax,\ymax);
				\draw[axis] (\xmin-0.2,0) -- (\xmax+0.2,0) node[axis label] {$x$};
				\draw[axis] (0,\ymin-0.2) -- (0,\ymax+0.2) node[axis label] {$y$};
				
				\foreach \j in {-120,...,\i} {
				  \pgfmathsetmacro{\a}{-0.025*\j}
				  \fill[point] ($(A)+\a*(V)$) circle (1pt);
				}
				
				\fill[current point] (A) circle (2pt) node[above] {$A$};
				
				\draw[vector] (A) -- ($(A)+\t*(V)$) node[vector label] {$t\cdot\vec{v}$};
				\draw[refvector] (1,1) -- (3,2) node[vector label] {$\vec{u}$};
				
				\node[info label] at (-2.8,2.5) 
                    {$ t=\pgfmathprintnumber[fixed,precision=3, zerofill]{\t} $};
			\end{tikzpicture} 
		}
    \end{animateinline}
  \end{frame}

  \begin{frame}{Parametrické vyjádření}
    \begin{alertblock}{Parametrické vyjádření přímky}
      Nechť je přímka $p$ dána bodem $A[a_{x};a_{y}]$ a vektorem $\vec{v}(v_{x};v
      _{y})$. Parametrickým vyjádřením přímky rozumíme všechny rovnice, které dokážou
      vyjádřit všechny body $P$, pro které platí $P \in p$:
      \[
        P=A+t\cdot \vec{v}\ \ \ \ t\in \mathbb{R}
      \]
      Někdy se uvádí dvě rovnice pro výpočet příslušných souřadnic bodu $P[x,y]$:
      \[
        x=a_{x}+t\cdot v_{x}
      \]
      \[
        y=a_{y}+t\cdot v_{y} \ \ \ \ t\in \mathbb{R}
      \]
    \end{alertblock}
  \end{frame}

  \section{Obecná rovnice přímky}
  \begin{frame}
    Tady rovnice
  \end{frame}

  \section{Řešení}
  \begin{frame}[label=Ř1]{Řešení cvičení - Posun a tvorba vektorů}
    \zpet{C1}
  \end{frame}

  \begin{frame}[label=Ř2]{Řešení cvičení - Střed úsečky a její velikost}
    \zpet{C2}
  \end{frame}

  \begin{frame}[label=Ř3]{Řešení cvičení - převod na goniometrický tvar}
    \zpet{C3}
  \end{frame}

  \begin{frame}[label=Ř4]{Řešení cvičení - převod na algebraický tvar}
    \zpet{C4}
  \end{frame}
  \begin{frame}[label=Ř5]{Řešení cvičení - násobení v GONIO a Moivreova věta}
    \zpet{C5}
  \end{frame}

  \begin{frame}[label=Ř6]{Řešení cvičení - lineární rovnice v \C}
    \zpet{C6}
  \end{frame}
\end{document}
